\documentclass{tufte-book}

% Font and Encoding
\usepackage{fontspec}
\setmainfont{Open Sans}[SmallCapsFont={}]

% Disable \nobibliography warning if no bibliography is used
\usepackage{etoolbox}
\AtBeginDocument{\let\nobibliography\relax}

% for the toc
\makeatletter
\renewcommand{\l@part}[2]{%
  \addvspace{1.5em}\noindent
  {\bfseries\Large\color{brightred}\MakeUppercase{Part #1}}\hfill
  {\bfseries\Large\color{brightred}#2}  % Correctly reference the page number
  \par\addvspace{1.0em}
}
\makeatother

\makeatletter
\renewcommand{\l@chapter}[2]{%
  \addvspace{1em}%
  {\bfseries\normalsize\color{brightred}#1}\hfill
  {\bfseries\normalsize\color{brightred}#2} % Ensures only the correct page number
  \par\addvspace{0.5em}
}
\makeatother

% for the parts titles
\usepackage{titlesec}  % Allows custom part title styling
\titleformat{\part}[display]
  {\Huge\bfseries\color{brightblue}}  % Ensures bright blue styling
  {\partname~\thepart}  % "Part X"
  {20pt}{}

% Updated Color Definitions (based on GIMP values)
\usepackage{xcolor}
\definecolor{brightred}{rgb}{0.976, 0.247, 0.118}  % Bright Red (Main Red)
\definecolor{darkred}{rgb}{0.631, 0.098, 0.047}    % Dark Red (Accent)
\definecolor{brightblue}{rgb}{0.0, 0.0, 1.0}       % Bright Blue (Main Blue)
\definecolor{darkblue}{rgb}{0.0, 0.0, 0.318}       % Dark Blue (Accent)

% Enable Page Numbers and Headers (Fix for tufte-book)
\usepackage{fancyhdr}
\pagestyle{fancy}

\fancyhf{}  % Clear header/footer
\fancyhead[LE]{\leftmark}  % Chapter on even pages
\fancyhead[RO]{\rightmark}  % Section on odd pages
\fancyfoot[C]{\thepage}  % Page number centered in footer
\renewcommand{\headrulewidth}{0.4pt}  % Header rule
\renewcommand{\footrulewidth}{0pt}    % No footer rule

% Colored boxes (less rounded)
\usepackage{tcolorbox}
\tcbuselibrary{skins}

% Colored boxes tyoes
\newtcolorbox{commandbox}[1][]{%
    colframe=brightblue, colback=brightblue!10, coltitle=white, % White title text
    boxrule=1pt, arc=1mm, left=5pt, right=5pt, top=5pt, bottom=5pt,
    fonttitle=\bfseries, title={#1}, colbacktitle=brightblue % Match background color
}

\newtcolorbox{configbox}[1][]{%
    colframe=brightred, colback=brightred!10, coltitle=white, % White title text
    boxrule=1pt, arc=1mm, left=5pt, right=5pt, top=5pt, bottom=5pt,
    fonttitle=\bfseries, title={#1}, colbacktitle=brightred % Match background color
}

% Code Formatting
\usepackage{listings}
\lstset{
    backgroundcolor=\color{darkblue!5},
    basicstyle=\ttfamily\footnotesize,
    frame=single,
    rulecolor=\color{darkblue},
    breaklines=true,
    keywordstyle=\color{brightblue}\bfseries,
    commentstyle=\color{darkred},
    stringstyle=\color{brightred}
}

% Define TOML Syntax
\lstdefinelanguage{TOML}{
    morekeywords={true,false},
    sensitive=true,
    morecomment=[l]{\#},  % TOML comments start with #
    morestring=[b]",      % Strings in double quotes
    morestring=[b]',      % Strings in single quotes
    keywordstyle=\color{brightblue}\bfseries,
    commentstyle=\color{darkred},
    stringstyle=\color{brightred},
    showspaces=false
}

% Define Bash Syntax
\lstdefinelanguage{Bash}{
    morekeywords={if, then, else, fi, for, while, do, done, exit, echo},
    sensitive=true,
    morecomment=[l]{\#},  % Comments start with #
    morestring=[b]",      % Double-quoted strings
    morestring=[b]',      % Single-quoted strings
    keywordstyle=\color{brightblue}\bfseries,
    commentstyle=\color{darkred},
    stringstyle=\color{brightred}
}

% Define Python Syntax
\lstdefinelanguage{Python}{
    morekeywords={def, return, if, elif, else, for, while, break, continue, import, from, as, class, try, except, with, yield, lambda, global, nonlocal, pass, raise, True, False, None},
    sensitive=true,
    morecomment=[l]{\#},  % Python comments start with #
    morestring=[b]",      % Double-quoted strings
    morestring=[b]',      % Single-quoted strings
    keywordstyle=\color{brightblue}\bfseries,
    commentstyle=\color{darkred},
    stringstyle=\color{brightred}
}

% Define Go Syntax
\lstdefinelanguage{Go}{
    morekeywords={package, import, func, return, var, const, type, struct, interface, map, range, if, else, for, switch, case, default, go, select, defer, break, continue, fallthrough},
    sensitive=true,
    morecomment=[l]{//},  % Single-line comment
    morecomment=[s]{/*}{*/}, % Multi-line comment
    morestring=[b]",      % Double-quoted strings
    morestring=[b]`,      % Backtick strings
    keywordstyle=\color{brightblue}\bfseries,
    commentstyle=\color{darkred},
    stringstyle=\color{brightred}
}

% Define R Syntax
\lstdefinelanguage{R}{
    morekeywords={if, else, repeat, while, function, for, in, next, break, TRUE, FALSE, NULL, NA, NaN, Inf},
    sensitive=true,
    morecomment=[l]{\#},  % R comments start with #
    morestring=[b]",      % Double-quoted strings
    morestring=[b]',      % Single-quoted strings
    keywordstyle=\color{brightblue}\bfseries,
    commentstyle=\color{darkred},
    stringstyle=\color{brightred}
}

% Define Julia Syntax
\lstdefinelanguage{Julia}{
    morekeywords={function, return, if, elseif, else, for, while, break, continue, import, using, struct, mutable, try, catch, finally, true, false, nothing},
    sensitive=true,
    morecomment=[l]{\#},  % Julia comments start with #
    morestring=[b]",      % Double-quoted strings
    morestring=[b]',      % Single-quoted strings
    keywordstyle=\color{brightblue}\bfseries,
    commentstyle=\color{darkred},
    stringstyle=\color{brightred}
}

% Colored Section Headings
\usepackage{sectsty}
\chapterfont{\color{brightred}}   % Chapter titles in bright red
\sectionfont{\color{darkred}}     % Section titles in dark red
\subsectionfont{\color{brightblue}} % Subsection titles in bright blue

% Hyperlinks
\usepackage{hyperref}
\hypersetup{
    colorlinks=true,
    linkcolor=brightred,
    urlcolor=brightblue,
    citecolor=darkred
}

% Define Custom Margin Note Commands
\newcommand{\tip}[1]{\marginnote{\textbf{\textcolor{brightblue}{Tip:}} #1}}
\newcommand{\warning}[1]{\marginnote{\textbf{\textcolor{brightred}{Warning:}} #1}}
\newcommand{\reminder}[1]{\marginnote{\textbf{\textcolor{darkred}{Reminder:}} #1}}
\newcommand{\note}[1]{\marginnote{\textbf{\textcolor{darkblue}{Note:}} #1}}

% Define a command to reference sections or chapters properly
\newcommand{\secref}[2]{\hyperref[#1]{\textbf{#2}}}

% for version tagging in document
\newcommand{\docversion}{v0.0.0
}

% Define title manually to bypass \maketitle issues
\newcommand{\customtitle}{
    \begin{titlepage}
        \centering
        \vspace*{3cm}
        \begin{figure}
            \centering
            \includegraphics{figures/prismAId_logo.png}
        \end{figure}
        {\Huge\bfseries\sffamily prism\textcolor{brightred}{A}\textcolor{brightblue}{I}d User Manual}\\[2cm]
        {\Large\textit{Riccardo Boero}}\\[0.3cm]
        \texttt{ribo@nilu.no}\\[1cm]
        {\large\today}\\[0.3cm]  % Email
        {\large Version: \docversion}\\[2cm]  % tagged version
        DOI: \href{https://doi.org/10.5281/zenodo.15025694}{10.5281/zenodo.15025694}  % DOI link

        % Creative Commons License
        \vfill
        \noindent
        \begin{minipage}{0.8\textwidth}
            \centering
            \textbf{Licensed under CC BY-SA 4.0} \\[0.5cm]
            \includegraphics[width=3cm]{figures/by-sa.png} \\[0.5cm]  % CC logo (adjust path)
            This work is licensed under the
            \href{https://creativecommons.org/licenses/by-sa/4.0/}{Creative Commons Attribution-ShareAlike 4.0 International License}.
        \end{minipage}
    \end{titlepage}
    \thispagestyle{plain}  % Force page number on the title page
}

\begin{document}
\lstset{showstringspaces=false}
\customtitle  % Use custom title instead of \maketitle

\tableofcontents  % Table of Contents after the title page
\newpage
%\pagestyle{plain}  % Ensure numbering continues

\chapter*{Foreword}
\addcontentsline{toc}{chapter}{Foreword}

The \texttt{prismAId} user manual is designed to help researchers, academics, and professionals leverage the power of the \texttt{prismAId} toolkit for conducting systematic reviews efficiently.\reminder{No coding skills are required to effectively conduct a systematic review with \texttt{prismAId}.} This document provides a structured approach to installing, configuring, and using the various components of \texttt{prismAId}, ensuring that users can quickly get started while also exploring advanced features when needed.

This manual is divided into five distinct parts, each catering to different user needs:

\begin{itemize}
    \item \textbf{Part 1: Introduction}\tip{The many open science advantages of \texttt{prismAId} are introduced and discussed in the \secref{chap:intro}{\textbf{Introduction}}.} -- If you are new to \texttt{prismAId}, we recommend starting with the chapter \secref{chap:intro}{\textbf{Introduction to prismAId}}. This section explains what \texttt{prismAId} is, the modular tools it provides, who can benefit from them, and why it is a valuable toolkit for systematic reviews.

    \item \textbf{Part 2: Getting Started}\note{\texttt{prismAId} is available on all platforms and operating systems. It can also be integrated programmatically with the most widely used scientific software.} -- If you need to install and configure \texttt{prismAId}, go directly to \secref{chap:install}{\textbf{Installation \& Setup}}. This section provides step-by-step installation instructions for Windows, Mac, and Linux.

    \item \textbf{Part 3: Conducting a Systematic Review}\tip{For a quick but complete walkthrough on using \texttt{prismAId} in a systematic review, see the \secref{chap:walkthrough}{\textbf{fifth}} chapter.} -- If your primary goal is to \textbf{learn how to conduct a systematic review with \texttt{prismAId}}, you can skip the installation and configuration sections and go directly to the chapter \secref{chap:walkthrough}{\textbf{Step-by-Step Guide to Conducting a Systematic Review}}. This section provides a hands-on walkthrough of:
    \begin{itemize}
        \item Setting up a project
        \item Acquiring and preparing literature using the Download and Convert tools
        \item Configuring and running systematic reviews using the Review tool
        \item Interpreting results
        \item Exporting findings
    \end{itemize}
    Before diving into this walkthrough, you may find it helpful to read \secref{chap:sysrev}{\textbf{Understanding Systematic Reviews}}, which provides background information on the methodology and best practices.

    \item \textbf{Part 4: Advanced Features} -- Users who want to explore advanced capabilities can refer to the chapters \secref{chap:advanced_ensemble}{\textbf{{chap:advanced_ensemble}}} and \secref{chap:advanced_features}{\textbf{Debugging, Cost Management \& Integration}}. These sections cover advanced options such as ensemble reviews, rate limiting, cost minimization, and integration with external tools and workflows.

    \item \textbf{Part 5: Troubleshooting \& FAQs} -- If you encounter issues, visit \secref{chap:troubleshooting}{\textbf{Troubleshooting Common Issues}}, which provides solutions to common errors for each of the toolkit's components, and \secref{chap:faq}{\textbf{Frequently Asked Questions}}, which addresses common concerns about usage and implementation.
\end{itemize}

\section*{How Should You Use This Manual?}
\addcontentsline{toc}{section}{How Should You Use This Manual?}

\begin{itemize}
    \item If you are \textbf{new to \texttt{prismAId}}, start with \secref{chap:intro}{\textbf{Introduction}}, then proceed to \secref{chap:install}{\textbf{Installation}}.
    \item If you \textbf{want to conduct a complete systematic review}, go directly to \textbf{Part 3}, particularly \secref{chap:walkthrough}{\textbf{Step-by-Step Guide to Conducting a Systematic Review}}, for a detailed walkthrough of the entire process.
    \item If you only need \textbf{specific tools} from the toolkit (Download, Convert, or Review), the respective chapters in \textbf{Part 3} detail how to use each component independently.
    \item If you need \textbf{advanced features or troubleshooting}, refer to \textbf{Parts 4 and 5} as needed.
\end{itemize}

Throughout this manual, useful commands are presented in blue boxes, as shown below:

\begin{commandbox}[Command: Example of a Binary Command]
To run a systematic review with a TOML configuration file:
\begin{lstlisting}[language=Bash]
# For Windows
./prismaid.exe -project your_project.toml

# For downloading papers from Zotero
./prismaid.exe -download-zotero zotero_config.toml

# For converting PDF files to text
./prismaid.exe -convert-pdf ./papers
\end{lstlisting}
\end{commandbox}

Tool configuration details are highlighted in red boxes, as shown below:

\begin{configbox}[Configuration: Example of a TOML Setting]
To configure an LLM for your review:
\begin{lstlisting}[language=TOML]
# Sample TOML configuration
[project.llm.1]
provider = "OpenAI"
api_key = ""
model = ""
temperature = 0.2
tpm_limit = 0
rpm_limit = 0
\end{lstlisting}
\end{configbox}


\section*{Final Notes}
\addcontentsline{toc}{section}{Final Notes}

We hope this manual serves as a comprehensive guide to making the most of the \texttt{prismAId} toolkit\warning{The online documentation of \texttt{prismAId} is available at \href{https://open-and-sustainable.github.io/prismaid/}{open-and-sustainable.github.io/prismaid/}.}. Whether you're a first-time user or an advanced researcher, this document is structured to help you get the information you need quickly and efficiently. The modular design of \texttt{prismAId} allows you to use only the components you need, making it adaptable to various research workflows and requirements.

This manual provides a complete overview of the features available in \texttt{prismAId} version ≥ 0.8.0. However, many aspects also apply to earlier versions that offer the same features.

For any additional questions, please refer to \secref{chap:faq}{the FAQ section} or contact our support team through the GitHub repository or Matrix Support Room. Happy reviewing!

\cleardoublepage
\markboth{}{}  % Clear any leftover chapter/section marks
\pagestyle{fancy}

% ==============================
% Part 1: Introduction
% ==============================
\part{Introduction}
\chapter{Introduction} \label{chap:intro}

Systematic reviews are at the core of evidence-based research. They help synthesize vast amounts of literature, ensuring that decisions in science, medicine, policy, and other fields are grounded in the best available evidence. However, conducting systematic reviews is a time-intensive and often overwhelming process. Screening literature, managing citations, extracting data, limiting biases and subjectivity, and ensuring methodological rigor and replicability all demand meticulous attention to detail.

\texttt{prismAId} was created to address these challenges. It is an open-source toolkit designed to assist researchers in conducting systematic reviews more efficiently using generative AI technology.\tip{Systematic reviews require structured workflows. \texttt{prismAId} provides modular tools to help enforce best practices at each stage of the process.} By streamlining key aspects of the review process, \texttt{prismAId} helps users maintain high standards of rigor and reproducibility while reducing the manual workload.

The toolkit is designed for a diverse range of users, from researchers conducting large-scale systematic reviews to students working on literature-based projects, requiring no coding skills to operate effectively. It provides structured workflows that align with best practices in systematic reviewing, ensuring that every step -- from acquiring literature to extracting and analyzing data -- is traceable and transparent.

\bigskip

\section{The prismAId Toolkit}

A key enhancement to \texttt{prismAId} is its modular design, which now offers four specialized tools to support different stages of the systematic review process:\note{Each tool can be used independently or as part of an integrated workflow. This modular design supports diverse research needs and practices.}

\begin{itemize}
    \item \textbf{Screening Tool}: Filters manuscripts after initial search but before full-text download, using deduplication, language detection, article type classification, and topic relevance scoring to identify papers for exclusion.

    \item \textbf{Download Tool}: Acquires papers directly from Zotero collections or from URL lists, streamlining the literature acquisition phase for manuscripts that passed screening.

    \item \textbf{Convert Tool}: Transforms documents from various formats (PDF, DOCX, HTML) into plain text that can be processed by large language models.

    \item \textbf{Review Tool}: Processes systematic literature reviews based on configurable protocols, extracting structured information from scientific papers.
\end{itemize}

\begin{figure}[h]
    \centering
    \includegraphics[width=\textwidth]{figures/prismaid_workflow_overview.png}
    \caption{The complete prismAId workflow: from literature search through screening, downloading, conversion, to final information extraction and analysis. Each tool can be used independently or as part of an integrated systematic review pipeline.}
    \label{fig:workflow_overview}
\end{figure}

This modular approach allows researchers to integrate \texttt{prismAId} into existing workflows more flexibly, using only the components needed for their specific research context.

\section{Open Science Foundation}

One of the key motivations behind \texttt{prismAId} is its commitment to Open Science. The toolkit is fully open source, meaning its development is transparent, and researchers can inspect, modify, and contribute to its functionality. This openness ensures that systematic reviews conducted with \texttt{prismAId} can be fully reproducible and that researchers can collaborate effectively without relying on proprietary or closed software ecosystems.\warning{As an open-source toolkit, \texttt{prismAId} does not include proprietary support. Users are encouraged to engage with the community through GitHub issues or the Matrix Support Room for troubleshooting and development.}

Moreover, the open-source nature of \texttt{prismAId} allows for continuous improvement through community-driven development. Users can adapt the tools to fit their needs, extend their capabilities, and integrate them with other research software. This flexibility is particularly beneficial in a rapidly evolving scientific landscape where reproducibility and adaptability are crucial.

\section{Multi-Platform Accessibility}

\texttt{prismAId} is accessible through multiple platforms and programming languages, offering flexibility based on user preference and technical requirements:

\begin{itemize}
    \item \textbf{Standalone Binaries}: For Windows, macOS, and Linux, requiring no coding skills
    \item \textbf{Go Package}: For integration in Go-based projects
    \item \textbf{Python Package}: For integration in Python scripts and Jupyter notebooks
    \item \textbf{R Package}: For use within R and RStudio environments
    \item \textbf{Julia Package}: For integration in Julia workflows
\end{itemize}

This multi-platform approach ensures that researchers can incorporate \texttt{prismAId} into their preferred computational environments without disrupting existing workflows.

\section{Guiding Principles}

Beyond technical efficiency, \texttt{prismAId} embodies key guiding principles that impact its users:

\begin{itemize}
    \item \textbf{Transparency}: The logic behind how studies are managed and processed is open for review. There are no hidden decision-making mechanisms.

    \item \textbf{Reproducibility}: Every action taken within \texttt{prismAId} can be logged and traced, allowing others to replicate results with comprehensive documentation of methodologies.

    \item \textbf{Flexibility}: Researchers can configure \texttt{prismAId} to meet the specific requirements of their systematic review protocols and adapt it to various research domains.

    \item \textbf{Accessibility}: Easy-to-use interfaces ensure that anyone can leverage advanced AI tools for literature reviews without specialized technical knowledge.

    \item \textbf{Efficiency}: The toolkit is optimized for handling large datasets with minimal setup, reducing the time from research to results.

    \item \textbf{No Vendor Lock-in}: Users are not dependent on a single commercial provider, ensuring that research workflows remain independent.

    \item \textbf{Community and Collaboration}:\reminder{Collaboration is a core feature. Consider contributing to \texttt{prismAId} if you have ideas for improvement!} As an open-source toolkit, \texttt{prismAId} fosters collaboration among researchers, developers, and systematic review specialists.
\end{itemize}

By embracing these principles, \texttt{prismAId} is more than just a toolkit -- it is part of a broader effort to improve the accessibility, efficiency, and reliability of systematic reviews. Whether used by a lone researcher or a large team, it provides the structure and support needed to navigate complex literature and synthesize high-quality evidence using the latest advancements in artificial intelligence.

As we move through this manual,\note{This manual follows a structured approach. Refer to the \textbf{How Should You Use This Manual?} section in the Foreword for guidance on navigating the content.} you will learn how to install, configure, and effectively use each component of the \texttt{prismAId} toolkit to conduct systematic reviews. From getting started with basic setup to leveraging advanced features for more complex reviews, this guide will provide all the necessary information to integrate \texttt{prismAId} into your research workflow.


% ==============================
% Part 2: Getting Started
% ==============================
\part{Getting Started}
\chapter{Installation \& Setup}
\label{chap:install}

This chapter explains how to install \texttt{prismAId}, covering system requirements, installation methods, and first-time setup for each of the toolkit's components.

\section{System Requirements}

Before installing \texttt{prismAId}, ensure your system meets the following requirements:\warning{Ensure that you have administrative privileges on your system before installing \texttt{prismAId}, especially on Windows and macOS.}

\begin{itemize}
    \item \textbf{Operating System}: Windows 10 or later, macOS 11 or later, or a Linux distribution (Arch, Debian, Fedora, Ubuntu, etc.).
    \item \textbf{Processor}: 64-bit CPU (Intel, AMD, or ARM64).
    \item \textbf{Memory}: At least 4GB RAM (8GB recommended).
    \item \textbf{Storage}: Minimum 500MB of free disk space.
    \item \textbf{Internet Connection}: Required for downloading packages and accessing LLM APIs.
    \item \textbf{API Keys}: To use the Review tool, you'll need an API key from at least one of the supported LLM providers (OpenAI, GoogleAI, Cohere, Anthropic, or DeepSeek).
\end{itemize}

\section{Installation Methods}

\texttt{prismAId} can be installed in multiple ways, depending on your preferences and workflow requirements.\warning{The capitalization of the letters in the prismAId tool follows the conventions of each programming language, ensuring consistency within each context but resulting in variations across packages and languages.}
 Choose the method that best suits your needs.

\subsection{Method 1: Standalone Binaries}

The simplest installation method is to download the pre-compiled binaries, which require no additional dependencies.\note{Binaries are standalone executables that work on Linux, macOS, and Windows, supporting both AMD64 and Arm64 architectures.}

\begin{enumerate}
    \item Navigate to the \href{https://github.com/open-and-sustainable/prismaid/releases}{GitHub Releases page}.
    \item Download the appropriate binary for your operating system.
    \item Place the file into a directory of your choice.
    \item On Linux and macOS, set the binaries as executable files.
\end{enumerate}

\subsubsection{Windows}
\reminder{On Windows, you may need to allow execution if prompted by security settings.}
\begin{commandbox}[Command: Running \texttt{prismAId} on Windows]
After downloading, navigate to the folder and run:
\begin{lstlisting}[language=Bash]
# Download papers from Zotero
./prismaid.exe -download-zotero zotero_config.toml

# Download papers from URL list
./prismaid.exe -download-URL paper_urls.txt

# Convert PDF files to text
./prismaid.exe -convert-pdf ./papers

# Run the Review tool with a project configuration
./prismaid.exe -project your_project.toml
\end{lstlisting}
\end{commandbox}


\subsubsection{macOS}
\warning{You may need to approve the application in \textbf{System Preferences} under \textbf{Security \& Privacy} before running it.}
\begin{commandbox}[Command: Running \texttt{prismAId} on macOS]
After downloading, give execution permissions and run:
\begin{lstlisting}[language=Bash]
chmod +x prismaid

# Download papers from Zotero
./prismaid -download-zotero zotero_config.toml

# Convert various file formats to text
./prismaid -convert-pdf ./papers
./prismaid -convert-docx ./papers
./prismaid -convert-html ./papers

# Run the Review tool with a project configuration
./prismaid -project your_project.toml
\end{lstlisting}
\end{commandbox}

\subsubsection{Linux}
\tip{If you receive a "Permission denied" error, try running \texttt{chmod +x} again or executing with \texttt{sudo}.}
\begin{commandbox}[Command: Running \texttt{prismAId} on Linux]
\begin{lstlisting}[language=Bash]
chmod +x prismaid

# Initialize a new project configuration interactively
./prismaid -init

# Download papers from a URL list
./prismaid -download-URL paper_urls.txt

# Run the Review tool with a project configuration
./prismaid -project your_project.toml
\end{lstlisting}
\end{commandbox}


\subsection{Method 2: Go Package}

If you are a Go developer, you can use \texttt{prismAId} as a Go package.\tip{The Go package offers the most comprehensive functionality, as it is the native implementation.}

\begin{commandbox}[Command: Installing the Go Package]
\begin{lstlisting}[language=Bash]
go get "github.com/open-and-sustainable/prismaid"
\end{lstlisting}
\end{commandbox}

\begin{commandbox}[Code: Using prismAId in Go]
\begin{lstlisting}[language=Go]
import "github.com/open-and-sustainable/prismaid"

// Download papers from Zotero
err := prismaid.DownloadZoteroPDFs(username, apiKey, collectionName, parentDir)

// Download from URL list
err := prismaid.DownloadURLList("path/to/urls.txt")

// Convert files to text
err := prismaid.Convert(inputDir, "pdf,docx,html")

// Run a systematic review
err := prismaid.Review(tomlConfigString)
\end{lstlisting}
\end{commandbox}

\subsection{Method 3: Python Package}

For Python users, \texttt{prismAId} is available as a package on PyPI.\note{The Python package works on Linux and Windows AMD64, and macOS Arm64.}

\begin{commandbox}[Command: Installing the Python Package]
\begin{lstlisting}[language=Bash]
pip install prismaid
\end{lstlisting}
\end{commandbox}

\begin{commandbox}[Code: Using prismAId in Python]
\begin{lstlisting}[language=Python]
import prismaid

# Download papers from Zotero
prismaid.download_zotero_pdfs("username", "api_key", "collection_name", "./papers")

# Download from URL list
prismaid.download_url_list("urls.txt")

# Convert files to text
prismaid.convert("./papers", "pdf,docx,html")

# Run a systematic review
with open("project.toml", "r") as file:
    toml_config = file.read()
prismaid.review(toml_config)
\end{lstlisting}
\end{commandbox}

\subsection{Method 4: R Package}

For R users, \texttt{prismAId} is available as a package on R-universe.\note{The R package works on Linux AMD64 and macOS Arm64.}

\begin{commandbox}[Command: Installing the R Package]
\begin{lstlisting}[language=R]
install.packages("prismaid", repos = c("https://open-and-sustainable.r-universe.dev", "https://cloud.r-project.org"))
\end{lstlisting}
\end{commandbox}

\begin{commandbox}[Code: Using prismAId in R]
\begin{lstlisting}[language=R]
library(prismaid)

# Download papers from Zotero
DownloadZoteroPDFs("username", "api_key", "collection_name", "./papers")

# Download from URL list
DownloadURLList("urls.txt")

# Convert files to text
Convert("./papers", "pdf,docx,html")

# Run a systematic review
toml_content <- paste(readLines("project.toml"), collapse = "\n")
RunReview(toml_content)  # Note the capitalization
\end{lstlisting}
\end{commandbox}

\subsection{Method 5: Julia Package}

For Julia users, \texttt{prismAId} is available as a package published in the Julia General Registry.\note{The Julia package works on Linux and Windows AMD64, and macOS Arm64.}

\begin{commandbox}[Command: Installing the Julia Package]
\begin{lstlisting}[language=Julia]
using Pkg
Pkg.add("PrismAId")
\end{lstlisting}
\end{commandbox}

\begin{commandbox}[Code: Using prismAId in Julia]
\begin{lstlisting}[language=Julia]
using PrismAId

# Download papers from Zotero
PrismAId.download_zotero_pdfs("username", "api_key", "collection_name", "./papers")

# Download from URL list
PrismAId.download_url_list("urls.txt")

# Convert files to text
PrismAId.convert("./papers", "pdf,docx,html")

# Run a systematic review
toml_config = read("project.toml", String)
PrismAId.run_review(toml_config)
\end{lstlisting}
\end{commandbox}

\section{Setting Up API Keys}

To use the Review tool of \texttt{prismAId}, you need to set up API keys for at least one of the supported LLM providers.\warning{Keep your API keys secure. Never share them in public repositories or unencrypted communications.}

\subsection{Obtaining API Keys}

\begin{itemize}
    \item \textbf{Anthropic}: Register at \href{https://www.anthropic.com/}{anthropic.com} and generate an API key.
    \item \textbf{Cohere}: Sign up at \href{https://cohere.com/}{cohere.com} and retrieve your API key from the dashboard.
    \item \textbf{DeepSeek}: Create an account at \href{https://platform.deepseek.com/}{platform.deepseek.com} and obtain an API key.
    \item \textbf{GoogleAI}: Create an account at \href{https://aistudio.google.com}{aistudio.google.com} and generate an API key.
    \item \textbf{OpenAI}: Register at \href{https://www.openai.com/}{openai.com} and obtain an API key from your account dashboard.
\end{itemize}

\subsection{Configuring API Keys}

There are two ways to configure your API keys:

\subsubsection{Environment Variables}

Set environment variables for your API keys:

\begin{commandbox}[Command: Setting API Key Environment Variables]
\begin{lstlisting}[language=Bash]
# For OpenAI
export OPENAI_API_KEY="your-openai-api-key"

# For GoogleAI
export GOOGLEAI_API_KEY="your-googleai-api-key"

# For Cohere
export COHERE_API_KEY="your-cohere-api-key"

# For Anthropic
export ANTHROPIC_API_KEY="your-anthropic-api-key"

# For DeepSeek
export DEEPSEEK_API_KEY="your-deepseek-api-key"
\end{lstlisting}
\end{commandbox}

\subsubsection{Configuration File}

Add your API keys directly in the TOML configuration file:\warning{Do not share or publish configuration files containing API keys. Always remove them before sharing.}

\begin{configbox}[Configuration: API Keys in TOML]
\begin{lstlisting}[language=TOML]
[project.llm.1]
provider = "OpenAI"
api_key = "your-openai-api-key"
model = "gpt-4o-mini"
temperature = 0.01
tpm_limit = 0
rpm_limit = 0
\end{lstlisting}
\end{configbox}

\tip{If both environment variables and configuration file entries are present, the configuration file values take priority.}

\section{Verifying the Installation}

Check if \texttt{prismAId} is correctly installed:\note{If you encounter issues, refer to the troubleshooting section or seek help from the community through GitHub issues or the Matrix Support Room.}

\begin{commandbox}[Command: Verifying Installation]
\begin{lstlisting}[language=Bash]
# For binaries
./prismaid --help

# For Python
python -c "import prismaid; print(prismaid.__version__)"

# For R
R -e "library(prismaid); cat(prismaid_version())"

# For Julia
julia -e "using PrismAId; println(PrismAId.version())"
\end{lstlisting}
\end{commandbox}

\section{Use in Jupyter Notebooks}

When using prismAId (versions ≤ 0.6.6) in Jupyter notebooks with Python, special handling may be required for interactive prompts:

\begin{commandbox}[Code: Using prismAId v < 0.6.6 in Jupyter Notebooks]
\begin{lstlisting}[language=Python]
import pty
import os
import time
import select

def run_review_with_auto_input(input_str):
    master, slave = pty.openpty()  # Create a pseudo-terminal

    pid = os.fork()
    if pid == 0:  # Child process
        os.dup2(slave, 0)  # Redirect stdin
        os.dup2(slave, 1)  # Redirect stdout
        os.dup2(slave, 2)  # Redirect stderr
        os.close(master)
        import prismaid
        prismaid.RunReviewPython(input_str.encode("utf-8"))
        os._exit(0)

    else:  # Parent process
        os.close(slave)
        try:
            while True:
                rlist, _, _ = select.select([master], [], [], 5)
                if master in rlist:
                    output = os.read(master, 1024).decode("utf-8", errors="ignore")
                    if not output:
                        break  # Process finished

                    print(output, end="")

                    if "Do you want to continue?" in output:
                        print("\n[SENDING INPUT: y]")
                        os.write(master, b"y\n")
                        time.sleep(1)
        finally:
            os.close(master)
            os.waitpid(pid, 0)  # Ensure the child process is cleaned up

# Load your review (TOML) configuration
with open("config.toml", "r") as file:
    input_str = file.read()

# Run the review function
run_review_with_auto_input(input_str)
\end{lstlisting}
\end{commandbox}

\section{Conclusion}

You have now installed \texttt{prismAId}. The toolkit offers five different installation methods, making it accessible across various platforms and programming environments. Each component (Download, Convert, and Review) can be used independently or as part of an integrated workflow.

In the next chapter, we'll explore how to configure your project in detail to leverage all of \texttt{prismAId}'s capabilities.


% ==============================
% Part 3: Conducting a Systematic Review
% ==============================
\part{Conducting a Systematic Review}
\chapter{Understanding Systematic Reviews} \label{chap:sysrev}
- What is a systematic review?  
- Key concepts \& methodology  
- Best practices  

\chapter{Step-by-Step Guide to Conducting a Systematic Review} \label{chap:walkthrough}
- Setting up a project  
- Importing and managing data  
- Running analyses  
- Interpreting results  
- Exporting findings  

% ==============================
% Part 4: Advanced Features
% ==============================
\part{Advanced Features}
\chapter[Ensemble Reviews]{Advanced Features: Ensemble Reviews} \label{chap:advanced_ensemble}

This chapter explores ensemble reviews, one of the most powerful advanced features of \texttt{prismAId}, enabling users to leverage multiple AI models simultaneously for enhanced reliability and insight.

\section{Understanding Ensemble Reviews}

Ensemble reviews represent a sophisticated approach to systematic literature analysis where multiple large language models (LLMs) are used to extract information from the same document set. This methodology parallels the traditional practice of having multiple human reviewers assess the same literature.

\subsection{Concept and Benefits}\tip{Ensemble reviews are particularly valuable when conducting high-stakes systematic reviews where accuracy and reliability are critical, such as in clinical guideline development or policy formation.}

\begin{itemize}
    \item \textbf{Definition}: An ensemble review utilizes two or more AI models, either from the same provider or across different providers, to independently extract information from each document.
    \item \textbf{Theoretical Foundation}: The approach is grounded in ensemble learning theory, where multiple algorithms collectively produce more accurate and robust results than individual algorithms.
\end{itemize}

\begin{configbox}[Configuration: Simple Ensemble with Two Models]
\begin{lstlisting}[language=TOML]
[project.llm.1]
provider = "OpenAI"
api_key = ""
model = "gpt-4o-mini"
temperature = 0.01
tpm_limit = 0
rpm_limit = 0

[project.llm.2]
provider = "OpenAI"
api_key = ""
model = "gpt-4o"
temperature = 0.01
tpm_limit = 0
rpm_limit = 0
\end{lstlisting}
\end{configbox}

\subsection{Key Advantages}

Ensemble reviews offer several significant benefits:

\begin{itemize}
    \item \textbf{Uncertainty Quantification}: Variations in model outputs highlight where information may be ambiguous or open to interpretation.
    \item \textbf{Increased Confidence}: Strong agreement across models suggests higher reliability of the extracted information.
    \item \textbf{Bias Reduction}: Different models may have different biases; using multiple models helps identify and mitigate these biases.
    \item \textbf{Robustness to Model Limitations}: Different models excel at different tasks; using multiple models compensates for individual weaknesses.
\end{itemize}

\section{Configuring Multi-Model Ensembles}

\texttt{prismAId} offers exceptional flexibility in configuring ensemble reviews, allowing you to combine models from the same provider, different providers, or a mix of both.

\subsection{Using Multiple Models from a Single Provider}

This approach allows you to compare different models with varying capabilities from the same provider:\tip{When using multiple models from the same provider, consider including a range of model sizes to balance cost, speed, and accuracy. For example, include both gpt-3.5-turbo and gpt-4o to compare a faster, more economical model with a more sophisticated one.}

\begin{configbox}[Configuration: Ensemble with Multiple OpenAI Models]
\begin{lstlisting}[language=TOML]
[project.llm.1]
provider = "OpenAI"
api_key = ""
model = "gpt-3.5-turbo"
temperature = 0.01
tpm_limit = 0
rpm_limit = 0

[project.llm.2]
provider = "OpenAI"
api_key = ""
model = "gpt-4o-mini"
temperature = 0.01
tpm_limit = 0
rpm_limit = 0

[project.llm.3]
provider = "OpenAI"
api_key = ""
model = "gpt-4o"
temperature = 0.01
tpm_limit = 0
rpm_limit = 0
\end{lstlisting}
\end{configbox}

\subsection{Cross-Provider Ensemble Configuration}

Using models from different providers can be particularly valuable, as these models may have been trained on different datasets and using different architectures:\warning{When using models from multiple providers, ensure you have valid API keys for each provider configured either in your TOML file or as environment variables. Missing or invalid keys will cause the review to fail for that provider.}

\begin{configbox}[Configuration: Cross-Provider Ensemble Example]
\begin{lstlisting}[language=TOML]
[project.llm.1]
provider = "OpenAI"
api_key = ""
model = "gpt-4o-mini"
temperature = 0.01
tpm_limit = 0
rpm_limit = 0

[project.llm.2]
provider = "GoogleAI"
api_key = ""
model = "gemini-1.5-flash"
temperature = 0.01
tpm_limit = 0
rpm_limit = 0

[project.llm.3]
provider = "Anthropic"
api_key = ""
model = "claude-3-haiku"
temperature = 0.01
tpm_limit = 0
rpm_limit = 0
\end{lstlisting}
\end{configbox}

\subsection{Comprehensive Five-Provider Ensemble}

For maximum diversity and robustness, you can configure an ensemble using models from all supported providers:\reminder{Pay close attention to the temperature settings in ensemble reviews. Lower temperatures (0.01-0.1) produce more deterministic outputs, making it easier to compare responses across models.}


\begin{configbox}[Configuration: Full Five-Provider Ensemble]
\begin{lstlisting}[language=TOML]
[project.llm.1]
provider = "OpenAI"
api_key = ""
model = "gpt-4o-mini"
temperature = 0.01
tpm_limit = 0
rpm_limit = 0

[project.llm.2]
provider = "GoogleAI"
api_key = ""
model = "gemini-1.5-flash"
temperature = 0.01
tpm_limit = 0
rpm_limit = 0

[project.llm.3]
provider = "Cohere"
api_key = ""
model = "command-r"
temperature = 0.01
tpm_limit = 0
rpm_limit = 0

[project.llm.4]
provider = "Anthropic"
api_key = ""
model = "claude-3-haiku"
temperature = 0.01
tpm_limit = 0
rpm_limit = 0

[project.llm.5]
provider = "DeepSeek"
api_key = ""
model = "deepseek-chat"
temperature = 0.01
tpm_limit = 0
rpm_limit = 0
\end{lstlisting}
\end{configbox}

\section{Analyzing Ensemble Results}

After running an ensemble review, you'll need specialized approaches to analyze and interpret the results from multiple models.

\subsection{Output Format and Structure}

When running an ensemble review, \texttt{prismAId} creates separate result files for each model used:

\begin{commandbox}[Example Output Files from Ensemble Review]
\begin{lstlisting}
results_OpenAI_gpt-4o-mini.csv
results_GoogleAI_gemini-1.5-flash.csv
results_Anthropic_claude-3-haiku.csv
\end{lstlisting}
\end{commandbox}

The individual model files contain the standard extraction results. Analyzing them, researchers whould prepare ensemble summary file containing consensus metrics and comparison data.

\subsection{Quantifying Model Agreement}

To analyze the level of agreement between models:\tip{For categorical variables, Cohen's kappa or Fleiss' kappa (for more than two models) can provide a more sophisticated measure of inter-model agreement than simple percentage agreement.}

\begin{commandbox}[R Code for Analyzing Model Agreement: Part 1]
\begin{lstlisting}[language=R]
library(tidyverse)

# Load results from different models
model1 <- read_csv("results_OpenAI_gpt-4o-mini.csv")
model2 <- read_csv("results_GoogleAI_gemini-1.5-flash.csv")
model3 <- read_csv("results_Anthropic_claude-3-haiku.csv")

# Join datasets
comparison <- model1 %>%
  select(filename, crop_type, climate_factor) %>%
  rename(crop_type_model1 = crop_type,
         climate_factor_model1 = climate_factor) %>%
  left_join(
    model2 %>%
      select(filename, crop_type, climate_factor) %>%
      rename(crop_type_model2 = crop_type,
             climate_factor_model2 = climate_factor),
    by = "filename"
  ) %>%
  left_join(
    model3 %>%
      select(filename, crop_type, climate_factor) %>%
      rename(crop_type_model3 = crop_type,
             climate_factor_model3 = climate_factor),
    by = "filename"
  )
\end{lstlisting}
\end{commandbox}
\begin{commandbox}[R Code for Analyzing Model Agreement: Part 2]
\begin{lstlisting}[language=R]
# Calculate agreement for categorical variables
comparison <- comparison %>%
  mutate(
    crop_type_agreement = case_when(
      crop_type_model1 == crop_type_model2 & crop_type_model2 == crop_type_model3 ~ "full_agreement",
      crop_type_model1 == crop_type_model2 | crop_type_model1 == crop_type_model3 | crop_type_model2 == crop_type_model3 ~ "partial_agreement",
      TRUE ~ "no_agreement"
    ),
    climate_factor_agreement = case_when(
      climate_factor_model1 == climate_factor_model2 & climate_factor_model2 == climate_factor_model3 ~ "full_agreement",
      climate_factor_model1 == climate_factor_model2 | climate_factor_model1 == climate_factor_model3 | climate_factor_model2 == climate_factor_model3 ~ "partial_agreement",
      TRUE ~ "no_agreement"
    )
  )

# Summarize agreement levels
agreement_summary <- comparison %>%
  summarize(
    crop_type_full_agreement = mean(crop_type_agreement == "full_agreement"),
    crop_type_partial_agreement = mean(crop_type_agreement == "partial_agreement"),
    crop_type_no_agreement = mean(crop_type_agreement == "no_agreement"),
    climate_factor_full_agreement = mean(climate_factor_agreement == "full_agreement"),
    climate_factor_partial_agreement = mean(climate_factor_agreement == "partial_agreement"),
    climate_factor_no_agreement = mean(climate_factor_agreement == "no_agreement")
  )

print(agreement_summary)
\end{lstlisting}
\end{commandbox}


\subsection{Visualizing Ensemble Results}

Visualizations can help identify patterns of agreement and disagreement:\note{For numerical data extraction (like percentages or measurements), consider using scatterplots to visualize the correlation between values extracted by different models, or boxplots to show the distribution of values across models.}

\begin{commandbox}[Python Code for Visualizing Ensemble Agreement: Part 1]
\begin{lstlisting}[language=Python]
import pandas as pd
import matplotlib.pyplot as plt
import seaborn as sns

# Load results from different models
model1 = pd.read_csv("results_OpenAI_gpt-4o-mini.csv")
model2 = pd.read_csv("results_GoogleAI_gemini-1.5-flash.csv")
model3 = pd.read_csv("results_Anthropic_claude-3-haiku.csv")

# Create comparison dataframe for a specific variable
comparison = pd.DataFrame({
    'filename': model1['filename'],
    'OpenAI': model1['crop_type'],
    'GoogleAI': model2['crop_type'],
    'Anthropic': model3['crop_type']
})

# Count agreements for each paper
comparison['agreement_count'] = comparison.apply(
    lambda row: len(set([row['OpenAI'], row['GoogleAI'], row['Anthropic']])),
    axis=1
)
\end{lstlisting}
\end{commandbox}
\begin{commandbox}[Python Code for Visualizing Ensemble Agreement: Part 2]
\begin{lstlisting}[language=Python]
# Create a heatmap of disagreements
# Convert to wide format with models as columns and papers as rows
heatmap_data = comparison.set_index('filename')[['OpenAI', 'GoogleAI', 'Anthropic']]

# Create categorical mapping for visualization
unique_values = sorted(list(set(
    heatmap_data['OpenAI'].tolist() +
    heatmap_data['GoogleAI'].tolist() +
    heatmap_data['Anthropic'].tolist()
)))
value_map = {value: i for i, value in enumerate(unique_values)}

# Apply mapping
for col in heatmap_data.columns:
    heatmap_data[col] = heatmap_data[col].map(value_map)

# Create heatmap
plt.figure(figsize=(12, len(heatmap_data) * 0.3))
sns.heatmap(heatmap_data, cmap='viridis', yticklabels=True)
plt.title('Model Agreement Visualization (Crop Type)')
plt.tight_layout()
plt.savefig('ensemble_agreement_heatmap.png', dpi=300)

# Create agreement distribution chart
agreement_counts = comparison['agreement_count'].value_counts().sort_index()
plt.figure(figsize=(10, 6))
agreement_counts.plot(kind='bar')
plt.title('Distribution of Model Agreement')
plt.xlabel('Number of Unique Answers (1 = Full Agreement)')
plt.ylabel('Number of Papers')
plt.tight_layout()
plt.savefig('agreement_distribution.png', dpi=300)
\end{lstlisting}
\end{commandbox}


\subsection{Decision Strategies for Conflicting Results}

When models disagree on extracted information, several strategies can be employed:\warning{Be cautious when using majority voting for numerical values, as this may lead to misleading results. For numerical data, consider using the median or mean, and examine the standard deviation to identify high-variance cases.}

\begin{itemize}
    \item \textbf{Majority Voting}: Use the value that the majority of models agree on.
    \item \textbf{Weighted Voting}: Assign higher weight to more capable models (e.g., GPT-4 over GPT-3.5).
    \item \textbf{Conservative Approach}: For papers with significant disagreement, flag for manual review.
    \item \textbf{Model-Specific Trust}: For certain types of information, one model may consistently outperform others.
\end{itemize}

\begin{commandbox}[Python Code for Majority Voting: Part 1]
\begin{lstlisting}[language=Python]
import pandas as pd
from collections import Counter

# Load results from different models
model1 = pd.read_csv("results_OpenAI_gpt-4o-mini.csv")
model2 = pd.read_csv("results_GoogleAI_gemini-1.5-flash.csv")
model3 = pd.read_csv("results_Anthropic_claude-3-haiku.csv")

# Prepare a consolidated results dataframe
consolidated = pd.DataFrame({'filename': model1['filename']})

# Apply majority voting for each extracted field
for field in ['crop_type', 'climate_factor', 'adaptation_strategies_discussed']:
    consolidated[field] = [
        Counter([m1, m2, m3]).most_common(1)[0][0]
        for m1, m2, m3 in zip(
            model1[field].fillna(''),
            model2[field].fillna(''),
            model3[field].fillna('')
        )
    ]

    # For fields where all models disagree, mark for review
    consolidated[f'{field}_review_needed'] = [
        len(set([m1, m2, m3])) == 3
        for m1, m2, m3 in zip(
            model1[field].fillna(''),
            model2[field].fillna(''),
            model3[field].fillna('')
        )
    ]
\end{lstlisting}
\end{commandbox}
\begin{commandbox}[Python Code for Majority Voting: Part 2]
\begin{lstlisting}[language=Python]
# For numerical fields, take the median
for field in ['yield_impact_percentage']:
    try:
        # Convert to numeric, errors='coerce' will convert non-numeric to NaN
        vals1 = pd.to_numeric(model1[field], errors='coerce')
        vals2 = pd.to_numeric(model2[field], errors='coerce')
        vals3 = pd.to_numeric(model3[field], errors='coerce')

        # Calculate median
        consolidated[field] = [
            pd.Series([v1, v2, v3]).median()
            for v1, v2, v3 in zip(vals1, vals2, vals3)
        ]

        # Calculate standard deviation to identify high variance
        consolidated[f'{field}_std'] = [
            pd.Series([v1, v2, v3]).std()
            for v1, v2, v3 in zip(vals1, vals2, vals3)
        ]

        # Flag for review if standard deviation is high
        consolidated[f'{field}_review_needed'] = consolidated[f'{field}_std'] > 2.0
    except:
        # Handle cases where the field might not exist or can't be processed
        pass

# Save consolidated results
consolidated.to_csv("results_consolidated.csv", index=False)
\end{lstlisting}
\end{commandbox}

\section{Case Studies in Ensemble Reviews}

To illustrate the practical application of ensemble reviews, let's examine some case scenarios and best practices.

\subsection{Case 1: Basic Verification Ensemble}

\begin{configbox}[Configuration: Verification Ensemble]
\begin{lstlisting}[language=TOML]
# A simple two-model ensemble using different providers
# for basic verification of extraction results

[project.llm.1]
provider = "OpenAI"
api_key = ""
model = "gpt-3.5-turbo"  # More economical option
temperature = 0.01
tpm_limit = 0
rpm_limit = 0

[project.llm.2]
provider = "Anthropic"
api_key = ""
model = "claude-3-haiku"  # Similar capability level
temperature = 0.01
tpm_limit = 0
rpm_limit = 0
\end{lstlisting}
\end{configbox}

\textbf{Use Case}: This configuration is suitable for routine systematic reviews where you want a basic cross-check between providers without significantly increasing costs. It uses comparably efficient models from different providers.

\textbf{Expected Outcome}: When these models agree, you can have higher confidence in the extraction. When they disagree, you might flag those specific data points for manual review.

\subsection{Case 2: Tiered Capability Ensemble}

\begin{configbox}[Configuration: Tiered Capability Ensemble]
\begin{lstlisting}[language=TOML]
# A three-model ensemble using progressively more capable models
# from the same provider for comparative analysis

[project.llm.1]
provider = "OpenAI"
api_key = ""
model = "gpt-3.5-turbo"  # Base level
temperature = 0.01
tpm_limit = 0
rpm_limit = 0

[project.llm.2]
provider = "OpenAI"
api_key = ""
model = "gpt-4o-mini"  # Mid level
temperature = 0.01
tpm_limit = 0
rpm_limit = 0

[project.llm.3]
provider = "OpenAI"
api_key = ""
model = "gpt-4o"  # Top level
temperature = 0.01
tpm_limit = 0
rpm_limit = 0
\end{lstlisting}
\end{configbox}

\textbf{Use Case}: This approach is useful for understanding how model capability affects extraction quality. It can help determine if investing in more advanced models provides meaningful improvements for your specific extraction tasks.

\textbf{Expected Outcome}: By comparing results across models of increasing capability, you can determine the minimum model level needed for reliable extraction, potentially optimizing future costs.

\subsection{Case 3: Comprehensive Cross-Provider Ensemble}

\begin{configbox}[Configuration: Comprehensive Cross-Provider Ensemble]
\begin{lstlisting}[language=TOML]
# A five-provider ensemble using top-tier models from each provider
# for maximum reliability in high-stakes reviews

[project.llm.1]
provider = "OpenAI"
api_key = ""
model = "gpt-4o"
temperature = 0.01
tpm_limit = 0
rpm_limit = 0

[project.llm.2]
provider = "GoogleAI"
api_key = ""
model = "gemini-1.5-pro"
temperature = 0.01
tpm_limit = 0
rpm_limit = 0

[project.llm.3]
provider = "Cohere"
api_key = ""
model = "command-r-plus"
temperature = 0.01
tpm_limit = 0
rpm_limit = 0

[project.llm.4]
provider = "Anthropic"
api_key = ""
model = "claude-3-opus"
temperature = 0.01
tpm_limit = 0
rpm_limit = 0

[project.llm.5]
provider = "DeepSeek"
api_key = ""
model = "deepseek-chat"
temperature = 0.01
tpm_limit = 0
rpm_limit = 0
\end{lstlisting}
\end{configbox}

\textbf{Use Case}: This configuration is ideal for high-stakes systematic reviews such as those informing clinical guidelines, policy decisions, or major meta-analyses where maximum reliability is essential.\tip{Start with a small test dataset when using comprehensive ensembles like this to estimate costs and review time before committing to your full dataset.}

\textbf{Expected Outcome}: This approach provides the highest confidence in results through diverse model architectures and training data. Points of unanimous agreement across all five providers suggest very high reliability.


\subsection{Best Practices for Ensemble Reviews}

Based on case studies and practical experience, here are recommended best practices:\reminder{When reporting results from ensemble reviews in publications, clearly state which models were used, how disagreements were resolved, and what level of agreement was observed for key findings.}

\begin{itemize}
    \item \textbf{Start Small}: Test your ensemble configuration on a small subset (5-10 papers) to identify any issues before running the full review.
    \item \textbf{Match Complexity to Need}: Use simpler ensembles for routine reviews and more comprehensive ensembles for high-stakes reviews.
    \item \textbf{Control Variables}: Keep temperature settings consistent across models to ensure fair comparisons.
    \item \textbf{Budget Appropriately}: Ensemble reviews will multiply your API costs by the number of models used; plan accordingly.
    \item \textbf{Document Model Versions}: Note the specific model versions used, as providers regularly update their models.
    \item \textbf{Analyze Disagreements}: Patterns in model disagreements often reveal ambiguities in your prompt or in the literature itself.
\end{itemize}

\chapter[Debugging, Costs \& Integration]{Advanced Features: Debugging, Cost Management \& Integration} \label{chap:advanced_features}

Building on our exploration of ensemble reviews, this chapter covers additional advanced features of \texttt{prismAId} that can enhance your systematic review process: debugging techniques, cost management strategies, advanced prompt engineering, and workflow integration.

\section{Debugging and Validation Techniques}

\texttt{prismAId} offers several advanced debugging features that can help identify and resolve issues in your systematic review process.

\subsection{Log Levels and Debugging Output}

The \texttt{log\_level} parameter in your configuration file controls the verbosity of debugging information:

\begin{configbox}[Configuration: Log Level Settings]
\begin{lstlisting}[language=TOML]
[project.configuration]
log_level = "high"  # Options: "low", "medium", "high"
\end{lstlisting}
\end{configbox}

Each level provides different information:\tip{When troubleshooting issues, start by setting the log level to "medium" to get more detailed console output. If the issue persists, switch to "high" for complete logging to file.}
\begin{itemize}
    \item \textbf{Low}: Minimal information, only essential status updates
    \item \textbf{Medium}: Detailed process information printed to the console
    \item \textbf{High}: Comprehensive logs saved to a file, including all API interactions
\end{itemize}

\begin{commandbox}[Command: Running with High Log Level]
\begin{lstlisting}[language=Bash]
# First set log_level to "high" in your config file
./prismaid -project your_project.toml

# Check the log file created in your project directory
cat your_project.log
\end{lstlisting}
\end{commandbox}

\subsection{Duplication for Consistency Testing}

The duplication feature runs the same review twice on identical inputs, allowing you to assess the consistency of model responses:\warning{Enabling duplication will double your API usage and costs. Use this feature selectively for testing or when consistency validation is critical.}

\begin{configbox}[Configuration: Enabling Duplication]
\begin{lstlisting}[language=TOML]
[project.configuration]
duplication = "yes"  # Options: "yes", "no"
\end{lstlisting}
\end{configbox}

The duplication process:
\begin{enumerate}
    \item Creates temporary copies of your input files
    \item Processes them as a separate batch
    \item Compares results for consistency
    \item Cleans up temporary files
\end{enumerate}

\subsection{Chain-of-Thought Justification}

The Chain-of-Thought (CoT) justification feature provides visibility into the model's reasoning process:

\begin{configbox}[Configuration: Enabling CoT Justification]
\begin{lstlisting}[language=TOML]
[project.configuration]
cot_justification = "yes"  # Options: "yes", "no"
\end{lstlisting}
\end{configbox}

When enabled, this creates a separate .txt file for each processed document containing:\tip{CoT justifications are invaluable for validating extraction quality, identifying potential misinterpretations, and understanding why models might be struggling with particular papers or information types.}
\begin{itemize}
    \item The model's reasoning for each extracted data point
    \item Key passages from the document that influenced the extraction
    \item Any uncertainty or alternative interpretations considered
\end{itemize}


\section{Cost and Rate Management}

\texttt{prismAId} provides several features to help manage costs and API rate limits.

\subsection{Token and Request Rate Limiting}
\note{Each provider has different rate limit structures. OpenAI limits by TPM (tokens per minute), Anthropic by RPM (requests per minute), and others may have tiered systems. Check the provider's documentation for current limits.}

\begin{configbox}[Configuration: Rate Limiting Settings]
\begin{lstlisting}[language=TOML]
[project.llm.1]
provider = "OpenAI"
api_key = ""
model = "gpt-4o-mini"
temperature = 0.01
tpm_limit = 100000  # Tokens per minute limit
rpm_limit = 60      # Requests per minute limit
\end{lstlisting}
\end{configbox}

These settings help you:
\begin{itemize}
    \item \textbf{Avoid Provider Rate Limiting}: Stay below provider-enforced limits\warning{These limits can change as providers update their services. Always check the most recent documentation from each provider for current rate limits.}
    \item \textbf{Balance Resource Usage}: Prevent spikes in API consumption
\end{itemize}

\begin{commandbox}[Example Provider Rate Limits]
\begin{lstlisting}
Provider  | Model            | Default TPM  | Default RPM
----------|------------------|--------------|------------
OpenAI    | gpt-3.5-turbo    | 60,000       | 3,500
OpenAI    | gpt-4o           | 10,000       | 500
Anthropic | claude-3-haiku   | N/A          | 5,000
GoogleAI  | gemini-1.5-flash | 4,000,000    | N/A
Cohere    | command-r        | 100,000      | 1,000
\end{lstlisting}
\end{commandbox}

Please note that not respecting rate limits may result in API errors and in not obtaining results for some manuscripts.

\subsection{Automatic Cost Minimization}

\texttt{prismAId} can automatically select the most cost-effective model by leaving the model field empty:\tip{Automatic model selection is particularly useful for heterogeneous document collections where some papers may be too large for smaller models.}

\begin{configbox}[Configuration: Cost Minimization]
\begin{lstlisting}[language=TOML]
[project.llm.1]
provider = "OpenAI"
api_key = ""
model = ""  # Empty value enables automatic model selection
temperature = 0.01
tpm_limit = 0
rpm_limit = 0
\end{lstlisting}
\end{configbox}

How automatic model selection works:
\begin{enumerate}
    \item \texttt{prismAId} calculates the token count for each document
    \item It determines which provider models can handle the document within token limits
    \item Among eligible models, it selects the most cost-efficient option
    \item Different documents may be processed by different models based on size
\end{enumerate}

\section{Advanced Prompt Engineering}

The quality of your systematic review results largely depends on how effectively you structure your prompts. Here are advanced strategies for optimizing extraction through prompt engineering.

\subsection{Optimizing Prompt Components}

Each component of the prompt structure serves a specific purpose and can be optimized:

\begin{configbox}[Configuration: Advanced Prompt Components]
\begin{lstlisting}[language=TOML]
[prompt]
persona = "You are an expert environmental scientist specializing in climate change impacts on agriculture, with experience conducting systematic reviews according to PRISMA guidelines."
task = "Analyze the scientific paper provided and extract specific information about climate change impacts on crop yields, methodology used, and adaptation strategies discussed."
expected_result = "You should output a JSON object containing values for each key specified in the review structure. Values must adhere exactly to the prescribed formats and vocabulary."
definitions = "'Crop yield' refers to harvestable production per unit of land area. 'Statistical significance' means p < 0.05 unless otherwise specified in the paper. 'Adaptation strategies' include any interventions meant to reduce negative climate impacts on agriculture."
example = "For instance, if the paper states 'wheat yields decreased by 5.2% (p < 0.01) per degree Celsius warming, with irrigation mitigating 40% of losses', you would extract: {\"crop_type\": \"wheat\", \"climate_factor\": \"temperature increase\", \"yield_impact\": \"-5.2% per °C\", \"statistical_significance\": \"yes\", \"adaptation_strategies_discussed\": \"yes\", \"adaptation_effectiveness\": \"40% loss reduction\"}"
failsafe = "If specific information is not clearly stated in the document, do not speculate or infer beyond reasonable scientific interpretation. Respond with an empty string for missing data. If the paper doesn't address the topic at all, use 'not addressed' for categorical fields."
\end{lstlisting}
\end{configbox}

Key strategies for each component:\note{Research suggests that more detailed and domain-specific prompts generally produce more accurate extractions, particularly for specialized scientific concepts.}

\begin{itemize}
    \item \textbf{Persona}: Include relevant expertise and methodological background to establish appropriate context.
    \item \textbf{Task}: Be specific about analytical depth and the nature of extraction required.
    \item \textbf{Expected Result}: Define the exact output format and emphasize adherence to specified value constraints.
    \item \textbf{Definitions}: Provide domain-specific terminology explanations, especially for potentially ambiguous concepts.
    \item \textbf{Example}: Show realistic examples that cover different data patterns and edge cases.
    \item \textbf{Failsafe}: Include clear guidelines for handling uncertainty, ambiguity, and missing information.
\end{itemize}


\subsection{Advanced Review Structure Patterns}

Beyond context setting above, you may implement more sophisticated patterns in your review structure:

\begin{configbox}[Configuration: Advanced Review Structure Patterns]
\begin{lstlisting}[language=TOML]
# Hierarchical extraction
[review.1]
key = "methodology_type"
values = ["experimental", "observational", "modeling", "review", "mixed", "other", ""]

[review.2]
key = "experimental_design"
values = ["randomized controlled trial", "non-randomized controlled trial", "before-after", "other", "not applicable", ""]

# Conditional extraction
[review.3]
key = "sample_size_reported"
values = ["yes", "no"]

[review.4]
key = "sample_size_value"
values = [""]

# Relational extraction
[review.5]
key = "crop_types_studied"
values = [""]

[review.6]
key = "primary_crop"
values = [""]

# Confidence assessment
[review.7]
key = "statistical_methods_quality"
values = ["high", "medium", "low", "not applicable", "cannot determine", ""]

\end{lstlisting}
\end{configbox}

Advanced patterns include:\tip{When using advanced patterns, ensure your prompt includes clear instructions on how these patterns relate to each other, particularly for conditional extractions.}

\begin{itemize}
    \item \textbf{Hierarchical Extraction}: Extract general categories first, then specific details based on those categories.
    \item \textbf{Conditional Extraction}: Use yes/no fields to establish presence, then extract specific values only if present.
    \item \textbf{Relational Extraction}: Capture relationships between different entities (e.g., primary crop among multiple crops studied).
    \item \textbf{Confidence Assessment}: Include quality assessments of methodological elements.
\end{itemize}


\subsection{Iterative Prompt Refinement}

The most effective prompts are developed through an iterative process:\warning{Changing prompts mid-review can introduce inconsistency. Once you begin your full review, avoid modifying the prompt unless absolutely necessary. If changes are required, consider re-processing previously analyzed papers.}

\begin{infobox}[Iterative Prompt Development Process]
\begin{lstlisting}
1. Start with a basic prompt and review structure
2. Test on 3-5 representative papers
3. Analyze results for:
   - Incorrect extractions
   - Missing information
   - Ambiguous responses
   - Inconsistent formatting
4. Identify patterns in errors or weaknesses
5. Refine prompt components and review structure
6. Test again on the same papers plus 2-3 new ones
7. Repeat until results match expectations
8. Document all prompt versions and their performance
\end{lstlisting}
\end{infobox}

For more complex extractions, consider dividing review items across multiple project configurations, allowing for more tailored context prompts with specific examples. While this approach can improve extraction accuracy, be aware that it increases costs as manuscripts must be processed multiple times.

\section{Workflow Integration and Automation}

\texttt{prismAId} can be integrated into broader research workflows and automated pipelines.\reminder{For very large reviews (hundreds or thousands of papers), consider implementing a database backend to store intermediate results and enable incremental processing.}

\subsection{Integration with Research Workflows}

\texttt{prismAId} can be integrated with broader research tools and workflows:

\begin{itemize}
    \item \textbf{Reference Management}: Integrate with Zotero, Mendeley, or EndNote for literature organization
    \item \textbf{Statistical Analysis}: Export results to R, SPSS, or specialized meta-analysis software
    \item \textbf{Collaborative Research}: Share configurations and results through version control systems
    \item \textbf{Publication Workflows}: Generate formatted tables and figures for manuscript inclusion
\end{itemize}

\section{Conclusion}

Advanced features of \texttt{prismAId} significantly enhance its capabilities beyond basic systematic reviews.\tip{Start with simpler configurations and gradually incorporate advanced features as you become more familiar with \texttt{prismAId}. Each feature adds power but also complexity, so build your expertise incrementally.} Ensemble reviews provide increased confidence and uncertainty quantification. Debugging features help identify and resolve issues in the extraction process. Cost management strategies optimize resource utilization, while advanced prompt engineering improves the quality and reliability of extractions. Finally, programmatic integration and automation enable seamless incorporation of \texttt{prismAId} into comprehensive research workflows.

By mastering these advanced features, researchers can conduct more sophisticated, reliable, and efficient systematic reviews that maintain the highest standards of methodological rigor while leveraging the power of artificial intelligence.


% ==============================
% Part 5: Troubleshooting \& FAQs
% ==============================
\part{Troubleshooting \& FAQs}
\chapter{Troubleshooting Common Issues} \label{chap:troubleshooting}

This chapter addresses common issues you might encounter when using \texttt{prismAId} and provides practical solutions to resolve them quickly.

\section{Installation Problems}

\subsection{Binary Installation Issues}

\begin{itemize}
    \item \textbf{Issue: "Permission denied" error when running the executable}

    \begin{commandbox}[Solution: Fix Permission Issues]
    \begin{lstlisting}[language=Bash]
    # For Linux/macOS
    chmod +x prismaid

    # For Windows (if using PowerShell with execution policy restrictions)
    Set-ExecutionPolicy -Scope Process -ExecutionPolicy Bypass
    \end{lstlisting}
    \end{commandbox}

    \item \textbf{Issue: Security warnings when running on macOS}
    \\
    On macOS, you may need to right-click the executable and select "Open" the first time you run it. After confirming once, you can run it normally.

    \item \textbf{Issue: Missing dependencies message}
    \\
    The standalone binaries should include all required dependencies. If you see dependency errors, you may have downloaded an incomplete file. Try downloading again or check for platform-specific versions.
\end{itemize}

\subsection{Package Installation Issues}
\warning{Package installation problems are often environment-specific. If you continue to experience issues, consider using the standalone binary instead, which doesn't require compiling or managing dependencies.}

\begin{itemize}
    \item \textbf{Issue: Python package installation fails}

    \begin{commandbox}[Solution: Python Installation Issues]
    \begin{lstlisting}[language=Bash]
    # Ensure pip is up-to-date
    python -m pip install --upgrade pip

    # Install with verbose output to see errors
    pip install prismaid --verbose

    # If C compiler errors occur, install build tools
    # For Windows:
    pip install --upgrade setuptools wheel
    # For Linux:
    sudo apt-get install build-essential python3-dev
    \end{lstlisting}
    \end{commandbox}

    \item \textbf{Issue: R package installation problems}

    \begin{commandbox}[Solution: R Installation Issues]
    \begin{lstlisting}[language=R]
    # Make sure you have the necessary system libraries
    # For Linux (Ubuntu/Debian):
    # sudo apt-get install r-base-dev libcurl4-openssl-dev libssl-dev

    # Try installing from source
    install.packages("prismaid", repos = "https://open-and-sustainable.r-universe.dev",
                    type = "source")

    # Check for error messages
    warnings()
    \end{lstlisting}
    \end{commandbox}

    \item \textbf{Issue: Julia package fails to build}

    \begin{commandbox}[Solution: Julia Installation Issues]
    \begin{lstlisting}[language=Julia]
    # Update registry
    using Pkg
    Pkg.Registry.update()

    # Try with build flag
    Pkg.add("PrismAId"; build=true)

    # Check Julia version (requires 1.6+)
    versioninfo()
    \end{lstlisting}
    \end{commandbox}
\end{itemize}

\section{Configuration Errors}

\subsection{TOML File Syntax Problems}
\tip{Use a TOML validator like \href{https://toml-online.com/}{toml-online.com} to check your configuration file for syntax errors before running \texttt{prismAId}.}

\begin{itemize}
    \item \textbf{Issue: "Failed to parse TOML" or similar error}

    \begin{configbox}[Common TOML Syntax Errors and Fixes]
    \begin{lstlisting}[language=TOML]
    # INCORRECT: Missing quotes around strings with special characters
    input_directory = C:\path\with\backslashes

    # CORRECT: Use quotes and forward slashes or escaped backslashes
    input_directory = "C:/path/with/forward/slashes"
    # OR
    input_directory = "C:\\path\\with\\escaped\\backslashes"

    # INCORRECT: Malformed arrays
    values = ["yes" "no"]

    # CORRECT: Use commas between array elements
    values = ["yes", "no"]

    # INCORRECT: Forgetting to close a section
    [project.llm.1
    provider = "OpenAI"

    # CORRECT: Include the closing bracket
    [project.llm.1]
    provider = "OpenAI"
    \end{lstlisting}
    \end{configbox}

    \item \textbf{Issue: Nested sections not recognized correctly}
    \\
    TOML sections must be properly nested. For example, `[project.llm.1]` is correct, but `[project][llm][1]` would be incorrect.
\end{itemize}

\subsection{Path-Related Problems}
\tip{Always use absolute paths in your configuration file to avoid path resolution issues. Relative paths can cause problems if you run \texttt{prismAId} from different directories.}

\begin{itemize}
    \item \textbf{Issue: "Directory not found" or "File not found" errors}

    \begin{commandbox}[Solution: Verify Paths]
    \begin{lstlisting}[language=Bash]
    # Check if the directory exists
    # Windows
    dir "C:\path\to\check"

    # Linux/macOS
    ls -la /path/to/check

    # Create directories if missing
    # Windows
    mkdir "C:\path\to\create"

    # Linux/macOS
    mkdir -p /path/to/create
    \end{lstlisting}
    \end{commandbox}

    \item \textbf{Issue: Path encoding problems with non-ASCII characters}
    \\
    If your paths contain non-ASCII characters (like accented letters or non-Latin scripts), make sure your configuration file is saved with UTF-8 encoding.
\end{itemize}

\subsection{API Key Configuration Issues}

\begin{itemize}
    \item \textbf{Issue: "API key not found" or authentication errors}

    \begin{commandbox}[Solution: Check API Key Setup]
    \begin{lstlisting}[language=Bash]
    # Check if environment variable is set
    # Windows
    echo %OPENAI_API_KEY%

    # Linux/macOS
    echo $OPENAI_API_KEY

    # Set environment variable if needed
    # Windows
    set OPENAI_API_KEY=your-api-key

    # Linux/macOS
    export OPENAI_API_KEY=your-api-key
    \end{lstlisting}
    \end{commandbox}
    You may also specify the API key in the review project configuration file.

    \item \textbf{Issue: API key works in environment but not in configuration}
    \\
    Check for extra spaces or invisible characters in your API key. Copy-pasting from certain sources can introduce these. Try re-typing the key manually if problems persist.
\end{itemize}

\section{Common Fixes for Each Tool}

\subsection{Download Tool Issues}

\begin{itemize}
    \item \textbf{Issue: Zotero downloads failing}
    \note{Zotero collection names are case-sensitive. Ensure your collection path exactly matches what appears in your Zotero library.}

    \begin{configbox}[Solution: Zotero Configuration Fixes]
    \begin{lstlisting}[language=TOML]
    # Correct format for Zotero collection path
    # INCORRECT
    group = Systematic Review/Climate Papers

    # CORRECT
    group = "Systematic Review/Climate Papers"

    # For group collections, use the group name
    group = "Group Name/Collection Name"
    \end{lstlisting}
    \end{configbox}

    \item \textbf{Issue: URL downloads timing out or failing}
    \\
    Some publishers block automated downloads. Consider using Zotero integration or manually downloading papers from repositories with strict access controls.
\end{itemize}

\subsection{Convert Tool Issues}

\begin{itemize}
    \item \textbf{Issue: PDF conversion producing empty or gibberish text}
    \\
    The Convert tool cannot extract text from scanned image PDFs or PDFs with strong copy protection. Use OCR software like Adobe Acrobat or Tesseract first, or manually transcribe critical sections.

    \item \textbf{Issue: PDF conversion missing text or with garbled formatting}
    \\
    PDFs with complex layouts (multiple columns, text boxes, etc.) may not convert perfectly. Consider manually editing the extracted text files to fix critical sections or remove unnecessary content.
\end{itemize}

\subsection{Review Tool Issues}

\begin{itemize}
    \item \textbf{Issue: Review stops after processing only a few files}
    \note{If hitting rate limits, configure appropriate `tpm\_limit` and `rpm\_limit` values in your configuration. This will automatically pace the requests to stay within provider limits.}

    \begin{commandbox}[Solution: Check Rate Limits and API Status]
    \begin{lstlisting}[language=Bash]
    # Set log level to high in your configuration
    # Then run the review to see detailed error messages
    ./prismaid -project your_project.toml

    # Check provider status pages
    # OpenAI: https://status.openai.com/
    # GoogleAI: https://status.cloud.google.com/
    \end{lstlisting}
    \end{commandbox}

    \item \textbf{Issue: Reviews producing unexpected or empty results}
    \reminder{Check the model's justifications and logs to understand why it's producing unexpected results. You may need to refine your prompt or review structure.}

    \begin{configbox}[Solution: Prompt Debugging]
    \begin{lstlisting}[language=TOML]
    # Enable Chain-of-Thought justification to see model reasoning
    [project.configuration]
    cot_justification = "yes"

    # Enable higher log level to see full prompts and responses
    log_level = "high"
    \end{lstlisting}
    \end{configbox}

\end{itemize}

\section{Performance and Resource Issues}

\begin{itemize}
    \item \textbf{Issue: Reviews running very slowly}
    \tip{For large reviews, consider processing subsets of your documents in parallel if your API plan allows.}

    \begin{configbox}[Solution: Performance Optimization]
    \begin{lstlisting}[language=TOML]
    # Balance rate limits for better performance
    [project.llm.1]
    provider = "OpenAI"
    api_key = ""
    model = "gpt-3.5-turbo"  # Use faster model for initial testing
    temperature = 0.01
    tpm_limit = 60000  # Set appropriate limits based on your plan
    rpm_limit = 3000
    \end{lstlisting}
    \end{configbox}

    \item \textbf{Issue: Out of memory errors with large documents}
    \\
    Very large documents may exceed the context window of LLMs or cause memory issues. Consider splitting extremely large papers into smaller sections or removing non-essential content before processing.
\end{itemize}

\section{Jupyter Notebook Issues}

\begin{itemize}
    \item \textbf{Issue: Python kernel crashing in Jupyter when running reviews}

    This issue may occur when the kernel cannot handle interactive prompts. The following workaround intercepts and automatically answers confirmation prompts.

    \begin{commandbox}[Solution: Jupyter Notebook Integration]
    \begin{lstlisting}[language=Python]
    # Use the workaround for versions <= 0.6.6
    import pty
    import os
    import time
    import select

    def run_review_with_auto_input(input_str):
        master, slave = pty.openpty()  # Create a pseudo-terminal

        pid = os.fork()
        if pid == 0:  # Child process
            os.dup2(slave, 0)  # Redirect stdin
            os.dup2(slave, 1)  # Redirect stdout
            os.dup2(slave, 2)  # Redirect stderr
            os.close(master)
            import prismaid
            prismaid.RunReviewPython(input_str.encode("utf-8"))
            os._exit(0)

        else:  # Parent process
            os.close(slave)
            try:
                while True:
                    rlist, _, _ = select.select([master], [], [], 5)
                    if master in rlist:
                        output = os.read(master, 1024).decode("utf-8", errors="ignore")
                        if not output:
                            break  # Process finished

                        print(output, end="")

                        if "Do you want to continue?" in output:
                            print("\n[SENDING INPUT: y]")
                            os.write(master, b"y\n")
                            time.sleep(1)
            finally:
                os.close(master)
                os.waitpid(pid, 0)  # Ensure the child process is cleaned up

    # Load your review (TOML) configuration
    with open("config.toml", "r") as file:
        input_str = file.read()

    # Run the review function
    run_review_with_auto_input(input_str)
    \end{lstlisting}
    \end{commandbox}
\end{itemize}

\section{Seeking Further Help}

If you continue to experience issues after trying the solutions in this chapter:\reminder{When seeking help, always include your \texttt{prismAId} version, operating system, error messages, and, if possible, a minimal reproducible example of your issue.}

\begin{itemize}
    \item \textbf{Check Documentation}: Review the online documentation at \href{https://open-and-sustainable.github.io/prismaid/}{open-and-sustainable.github.io/prismaid/}

    \item \textbf{GitHub Issues}: Search existing issues or create a new one at \href{https://github.com/open-and-sustainable/prismaid/issues}{github.com/open-and-sustainable/prismaid/issues}

    \item \textbf{Matrix Support Room}: Join the \href{https://matrix.to/#/#prismAId-support:matrix.org}{prismAId Support Room} to discuss your issue with the community
\end{itemize}

\begin{commandbox}[Gathering System Information for Support]
\begin{lstlisting}[language=Bash]
# Get prismAId version
./prismaid --version

# For Python package
python -c "import prismaid; print(prismaid.__version__)"

# System information
# Windows
systeminfo | findstr /B /C:"OS Name" /C:"OS Version"

# Linux
uname -a
cat /etc/os-release

# macOS
sw_vers
\end{lstlisting}
\end{commandbox}

Remember that \texttt{prismAId} is an open-source project. Detailed bug reports help improve the software for everyone, and community contributions to fix issues are always welcome.

\chapter{Frequently Asked Questions} \label{chap:faq}

This chapter addresses common questions about \texttt{prismAId}, its capabilities, integration with other tools, and performance optimization.

\section{General Questions}

\subsection{What types of systematic reviews is prismAId suitable for?}

\texttt{prismAId} is designed to handle a wide range of systematic review types across various disciplines. It works well for:\note{While prismAId can extract information from any text-based literature, it performs best when the information being sought is explicitly stated in the text rather than requiring complex inference or domain-specific interpretation.}

\begin{itemize}
    \item Literature reviews in scientific fields
    \item Scoping reviews to map available evidence
    \item Meta-analyses when combined with statistical tools
    \item Policy reviews and evidence summaries
    \item Qualitative evidence syntheses
\end{itemize}

\subsection{Can I use prismAId with non-English literature?}

Yes, but with some limitations.\tip{For multilingual reviews, consider testing the extraction quality on a sample of non-English papers before proceeding with the full review. You may need to adjust your prompts to account for language-specific considerations.} \texttt{prismAId}'s underlying LLMs support multiple languages with varying degrees of proficiency. However:

\begin{itemize}
    \item English literature typically yields the most reliable results
    \item Major European languages (French, German, Spanish) generally work well
    \item Other languages may have reduced extraction accuracy
    \item The Convert tool may struggle with non-Latin character sets
\end{itemize}

\subsection{Is prismAId compliant with systematic review reporting guidelines?}

\texttt{prismAId} is designed to support guideline-compliant reviews but doesn't enforce reporting requirements itself:

\begin{itemize}
    \item It facilitates PRISMA-compliant reviews by enabling structured data extraction
    \item The methodology is transparent and reproducible, supporting various reporting frameworks
    \item You must still ensure your final report addresses all required elements from guidelines like PRISMA, MOOSE, or ENTREQ
\end{itemize}

\begin{commandbox}[Example Citation for prismAId in Methods Section]
\begin{lstlisting}
In our systematic review, data extraction was performed using prismAId
version X.X.X (Boero, 2024), an open-source tool that uses large
language models to extract structured information from scientific
literature. We used [MODEL NAMES] with a temperature setting of
[VALUE] to extract data according to our pre-defined protocol.
To ensure extraction quality, we manually validated [XX%] of
extracted data points, finding [XX%] agreement between AI-extracted
and manually verified information.

Boero, R. (2024). prismAId - Open Science AI Tools for Systematic,
Protocol-Based Literature Reviews. Zenodo.
https://doi.org/10.5281/zenodo.11210796
\end{lstlisting}
\end{commandbox}

\section{Integration with Other Software}

\subsection{Can I use prismAId with Zotero/Mendeley/EndNote?}

Yes, \texttt{prismAId} offers integration options with reference managers:\reminder{When using the Zotero integration, ensure papers have attached PDFs in your Zotero library. prismAId cannot download papers that don't have attachments.}

\begin{itemize}
    \item \textbf{Zotero}: Direct integration via the Download tool, which can access papers from your Zotero collections using the API
    \item \textbf{Mendeley/EndNote}: No direct API integration, but you can export papers to a folder and then use \texttt{prismAId} to process them
\end{itemize}

\subsection{Can I use prismAId with statistical software for meta-analysis?}

Yes, \texttt{prismAId} works well with statistical tools for meta-analysis:

\begin{itemize}
    \item Export results as CSV or JSON
    \item Import into R (with packages like \texttt{meta} or \texttt{metafor})
    \item Import into STATA, SPSS, or specialized meta-analysis software
    \item Use with Python's \texttt{metapsy} or similar packages
\end{itemize}

\begin{commandbox}[R Code for Importing prismAId Results for Meta-Analysis]
\begin{lstlisting}[language=R]
# Load libraries
library(readr)
library(dplyr)
library(meta)

# Import prismAId results
results <- read_csv("results.csv")

# Prepare data for meta-analysis
# Assuming prismAId extracted effect sizes, sample sizes, and p-values
meta_data <- results %>%
  # Clean and convert text to numeric values
  mutate(
    effect_size = as.numeric(gsub("[^0-9.-]", "", effect_size)),
    std_error = as.numeric(gsub("[^0-9.]", "", std_error)),
    sample_size = as.numeric(gsub("[^0-9]", "", sample_size))
  ) %>%
  # Remove rows with missing data
  filter(!is.na(effect_size) & !is.na(std_error))

# Run meta-analysis
if(nrow(meta_data) >= 3) {
  ma <- metagen(
    TE = meta_data$effect_size,
    seTE = meta_data$std_error,
    studlab = meta_data$filename,
    sm = "SMD"
  )

  # Print results
  summary(ma)

  # Create forest plot
  forest(ma)
}
\end{lstlisting}
\end{commandbox}

\subsection{Can I use prismAId with qualitative analysis software?}

Yes, \texttt{prismAId} can be used alongside qualitative analysis tools:\tip{For qualitative reviews, consider using more open-ended extraction fields in your review structure and then conducting more detailed thematic analysis in specialized software.}

\begin{itemize}
    \item Export results to CSV/JSON and import into NVivo, ATLAS.ti, or MAXQDA
    \item Use \texttt{prismAId} for initial coding and then refine in specialized software
    \item Combine AI-assisted extraction with manual qualitative analysis
\end{itemize}

\subsection{Can I integrate prismAId into automated research workflows?}

Yes, \texttt{prismAId} can be integrated into automated research workflows:

\begin{itemize}
    \item Use the API from Go, Python, R, or Julia in scripted workflows
    \item Incorporate into continuous integration pipelines
    \item Schedule batch processing of papers
    \item Chain with other research tools using scripts
\end{itemize}

\section{Performance Optimization}

\subsection{How can I optimize prismAId for speed?}

To improve processing speed:\warning{Faster models may have lower accuracy for complex extraction tasks. Always validate a sample of results when optimizing for speed.}

\begin{itemize}
    \item \textbf{Use faster models}: Models like \texttt{gpt-3.5-turbo}, \texttt{claude-3-haiku}, or \texttt{gemini-1.5-flash} process text more quickly than their larger counterparts
    \item \textbf{Reduce text length}: Remove non-essential sections (references, acknowledgments) from papers
    \item \textbf{Batch processing}: Process papers in parallel using multiple instances, an din particular if using multiple models and from different providers
    \item \textbf{Optimize rate limits}: Set appropriate \texttt{tpm\_limit} and \texttt{rpm\_limit} values based on your API plan
\end{itemize}

\begin{configbox}[Configuration: Speed Optimization]
\begin{lstlisting}[language=TOML]
[project.llm.1]
provider = "OpenAI"
api_key = ""
model = "gpt-3.5-turbo"  # Faster model
temperature = 0.01
tpm_limit = 60000  # Adjust based on your tier
rpm_limit = 3000   # Adjust based on your tier
\end{lstlisting}
\end{configbox}

\subsection{How can I optimize prismAId for accuracy?}

To improve extraction accuracy:

\begin{itemize}
    \item \textbf{Use more capable models}: Models like \texttt{gpt-4o}, \texttt{claude-3-opus}, or \texttt{gemini-1.5-pro} typically provide more accurate extractions
    \item \textbf{Refine prompts}: Create clear, specific prompts with good examples
    \item \textbf{Use ensemble reviews}: Combine multiple models to improve reliability
    \item \textbf{Lower temperature}: Set temperature to 0.01-0.1 for more deterministic responses
    \item \textbf{Break down complex extractions}: Use multiple simpler reviews instead of one complex review
\end{itemize}

\begin{configbox}[Configuration: Accuracy Optimization]
\begin{lstlisting}[language=TOML]
[project.llm.1]
provider = "OpenAI"
api_key = ""
model = "gpt-4o"  # More capable model
temperature = 0.01  # Minimal randomness
tpm_limit = 10000
rpm_limit = 500

[prompt]
# More detailed prompt components
persona = "You are an expert scientist with extensive experience conducting systematic reviews following PRISMA guidelines. You have specific expertise in identifying methodological details and extracting precise information from scientific papers."
# ... other detailed prompt components
\end{lstlisting}
\end{configbox}

\subsection{How can I optimize prismAId for cost efficiency?}

To minimize costs:\tip{For cost-sensitive projects, run a small test first with 5-10 papers and examine the log output (set `log\_level = "medium"`) to see token usage and estimated costs per paper.}

\begin{itemize}
    \item \textbf{Use automatic model selection}: Leave the model field empty (`model = ""`) to let \texttt{prismAId} choose the most cost-effective model
    \item \textbf{Trim input texts}: Remove unnecessary sections like references, acknowledgments, and front matter
    \item \textbf{Focus reviews}: Extract only the specific information you need rather than general information
    \item \textbf{Test on samples}: Validate your approach on a small sample before processing the entire dataset
\end{itemize}

\begin{configbox}[Configuration: Cost Optimization]
\begin{lstlisting}[language=TOML]
[project.llm.1]
provider = "OpenAI"
api_key = ""
model = ""  # Empty for automatic selection
temperature = 0.01
tpm_limit = 0
rpm_limit = 0
\end{lstlisting}
\end{configbox}

\section{Methodological Questions}

\subsection{How accurate is information extraction with prismAId?}

The accuracy of \texttt{prismAId} extractions depends on several factors:\reminder{Always validate a subset of extractions manually to assess accuracy for your specific review task. Consider using ensemble reviews for high-stakes systematic reviews.}

\begin{itemize}
    \item \textbf{Prompt quality}: Well-crafted prompts improve accuracy
    \item \textbf{Information type}: Explicit data is extracted more accurately than implicit information
    \item \textbf{Model capability}: More advanced models generally provide higher accuracy
    \item \textbf{Document quality}: Clean, well-structured texts yield better results
\end{itemize}

\subsection{How should I validate prismAId results?}

Best practices for validation include:

\begin{itemize}
    \item Manually check a random sample (10-20\%) of papers
    \item Use dual extraction (AI + human) for critical information
    \item Run ensemble reviews with multiple models and examine disagreements
    \item Calculate inter-rater reliability between AI and human extractions
    \item Document validation process and results in your methods section
\end{itemize}

\subsection{Can prismAId completely replace manual review?}

No, \texttt{prismAId} should be viewed as an assistive tool rather than a complete replacement for human judgment:\warning{Complete automation of systematic reviews without human supervision is not recommended, particularly for reviews that inform clinical practice, policy decisions, or other high-stakes applications.}

\begin{itemize}
    \item It excels at reducing the time-intensive and subjective aspects of systematic reviews
    \item Human oversight remains essential for quality control and interpretation
    \item Domain expertise is still needed to contextualize findings and draw conclusions
    \item Complex or nuanced information may require human validation
\end{itemize}

\section{Advanced Usage Questions}

\subsection{Can I use prismAId with custom or proprietary LLMs?}\tip{As an open-source project, prismAId can be extended to support additional LLM providers. If you have programming experience, you could fork the repository and add support for your preferred models by extending the interface code.}

Currently, \texttt{prismAId} supports five major commercial LLM providers (OpenAI, GoogleAI, Anthropic, Cohere, DeepSeek). Support for custom or self-hosted models is not built into the main distribution.

\subsection{Can I use prismAId for sensitive or confidential data?}

You should exercise caution when using \texttt{prismAId} with sensitive data:\warning{When using commercial LLM APIs, the text of your documents is sent to third-party servers. While providers typically have privacy safeguards, you should never process confidential or personally identifiable information without ensuring compliance with relevant regulations and obtaining appropriate permissions.}

\begin{itemize}
    \item All text sent to LLM providers may be processed on their servers
    \item Check each provider's privacy policy and data usage terms
    \item Consider data residency and compliance requirements (GDPR, HIPAA, etc.)
    \item For highly sensitive content, consider redacting identifying information before processing
\end{itemize}

\subsection{How can I contribute to prismAId development?}

As an open-source project, \texttt{prismAId} welcomes contributions:\reminder{Before contributing code, familiarize yourself with the project's contribution guidelines and code of conduct in the repository.}

\begin{itemize}
    \item Report bugs or request features through GitHub issues
    \item Contribute code improvements via pull requests
    \item Help improve documentation or examples
    \item Share validation studies or use cases
\end{itemize}

\begin{commandbox}[Contributing to prismAId]
\begin{lstlisting}[language=Bash]
# Fork the repository on GitHub
# Clone your fork
git clone https://github.com/your-username/prismaid.git

# Create a feature branch
git checkout -b feature/your-feature-name

# Make your changes
# Test your changes

# Commit with a descriptive message
git commit -m "Add feature: description of your changes"

# Push to your fork
git push origin feature/your-feature-name

# Create a pull request on GitHub
\end{lstlisting}
\end{commandbox}

\subsection{Will my API usage for prismAId affect my other projects using the same API keys?}

Yes, any \texttt{prismAId} usage will count toward your overall API usage limits with each provider:\tip{For organizations with multiple projects using the same LLM providers, consider creating separate API keys for different projects to track usage more effectively and implement budget controls.}

\begin{itemize}
    \item API calls made by \texttt{prismAId} will appear in your provider dashboard
    \item Usage counts toward any quotas or rate limits on your account
    \item Billing is based on total token usage across all applications using your key
\end{itemize}

\section{Future Development Questions}

\subsection{What new features are planned for prismAId?}

While specific roadmap details may change, planned enhancements include:

\begin{itemize}
    \item Support for additional LLM providers and models
    \item Improved OCR integration for image-based PDFs
    \item More sophisticated ensemble methods
\end{itemize}

\note{As an open-source project, prismAId's development direction is influenced by community needs and contributions. You can suggest features or vote on existing proposals through the GitHub repository.}

\subsection{How can I stay updated on prismAId developments?}

To stay informed about new releases and features:\reminder{Major version updates may include changes to configuration structure or API methods. Always review the changelog before upgrading.}

\begin{itemize}
    \item Follow the GitHub repository
    \item Join the Matrix Announcements Room
    \item Subscribe to release notifications
    \item Check the documentation website periodically
\end{itemize}

\begin{commandbox}[Following prismAId Updates]
\begin{lstlisting}
GitHub Repository: https://github.com/open-and-sustainable/prismaid
Matrix Announcements: https://matrix.to/#/#prismAId-announcements:matrix.org
Documentation: https://open-and-sustainable.github.io/prismaid/
\end{lstlisting}
\end{commandbox}

\section{Conclusion}

This FAQ covers common questions about \texttt{prismAId}, but the field of AI-assisted systematic reviews is rapidly evolving. For the most current information, consult the online documentation, join the community support channels, or review recent publications about \texttt{prismAId} methodology.

If your question isn't addressed here, consider posting in the Matrix Support Room or opening a discussion on the GitHub repository to engage with the community of users and developers.


\end{document}
